\paragraph{Theorem D.2.38 (Proved: the mixed ten- and fourteen-torsion smooth types).}
Take one of the twelve minimal patterns in the following table,
with the primitive Mumford filling at every multiplicative position,
the indicated additive choice, and no additional regular logarithmic
transformation. Retain the ordinary clutch $e\gamma$ and allow every
finite multiplicity and affine twist for which the output $X$ is
simply connected. Zero positions use the zero-based source indexing.
Let $Z_{p,\epsilon}$ denote the oriented smooth type of D.2.36 with
$H_2=\mathbb Z/p$, $b_3=0$, $\mu=0$, and refined Pontryagin
orbit $\mathfrak p=\epsilon$, for $\epsilon=0$ or $1$ when that
type occurs there. Then
\[
                  X\cong M_{\mathrm{conic}}\# Z_{p,\epsilon}
                     \quad\hbox{orientation-preservingly},
                                                        \tag{D.MP1}
\]
where the parameters are
\[
\begin{array}{c|r|r|l|l|r|r}
\text{OS/profile}&a&b&\text{zero positions}&\text{additive choice}&p&\epsilon\\\hline
56&1&10&1&\text{none}&5&0\\
56&3&10&3,4&\text{none}&5&1\\
56/II&1&10&3&II\text{ original or log}&5&0\\
56/IV&1&30&3&IV\text{ log}&5&0\\
56/IV&3&10&3&IV\text{ original}&5&1\\
47&1&14&2,3,4&\text{none}&7&1\\
47/II&1&14&2&II\text{ original or log}&7&1.
\end{array}                                               \tag{D.MP2}
\]
There are seven $c=10$ patterns and five $c=14$ patterns on ten
distinct marked initial rows. Thus they give exactly two oriented
smooth types for $c=10$ and one for $c=14$.

\begin{proof}
\emph{The actual classes and local extensions.}
D.2.26 gives $H_2(X;\mathbb Z)=H_3(X;\mathbb Z)=\mathbb Z/(2p)$,
with $b_3=0$. In particular $H^2(X;\mathbb Z)=0$, so the complex
first Chern class is zero and $X$ is spin. Use the lattice basis
$(u,w,\delta,\gamma)$ and the primitive alternating form
\[
             \Omega=a\,u^*w^*+b\,\gamma^*\delta^*.
                                                        \tag{D.MP3}
\]
Here $a$ is odd, $b/2$ is odd, and $p\mid b$. The coefficient
$b/q_i$ in every finite-cap normal relation is divisible by $2p$:
the only possibilities are $b/q_i=10$ or $14$, with the $IV$
logarithmic cap using $b/q_i=30/3=10$. The other coefficient is
a multiple of $b$. The allowed original additive caps have zero
fibre coefficient, and the ordinary clutch has coefficient $be$.
Thus the pure fibre functional, zero on the rigidified sweeps,
descends both modulo two and modulo $p$ through the actual
attachments for every allowed finite order and twist.

Let $W$ retain all multiplicative Mumford caps and omit the additive
caps and one ordinary disc, and let $V$ be its regular part. The
exact marked angular lattices in all ten rows have Hermite basis
$(u,w,\delta)$; so do the full images of $A_i-I$. Hence
$\pi_1(W)=\mathbb Z_\gamma\times F_{r-1}$, where $r$ is the number
of unfilled discs. The boundary-to-cap degree-two maps are onto.
As in the rank calculation (D.C26), the independent peripheral
base classes, with the ordinary disc still unfilled, give
$H_2(V)\twoheadrightarrow H_2(W)$. Both first homology groups are
free. Therefore restriction from $W$ to $V$ is injective in degree
two with integral, mod-two, or mod-$p$ coefficients. This proves
the needed comparison with the globally descended classes.

Use the unsquared symmetric factors (D.C27) for $\Omega$.
All their exponents are integral since $b$ is even. Their
semicharacter $q=a(u+w+uw)$ is preserved by every actual
monodromy, as the exact certificate checks on all sixteen
vectors at each site. Thus the rigidified factors patch.
Their comparison with the marked prequantum family has globally
two-torsion Chern class; $H^2(V;\mathbb Z)$ is free, so that class
is zero. The resulting line has exactly (D.MP3) on $V$,
including its zero swept values.

For completeness the local extension uses the raw Tate periods
\[
 H=(n,ar),\qquad F=(-br,\rho).
\]
At the rank-two zero they have determinant one, and
\[
 \phi(x,y)=\frac{-a\rho x^2-2ab rxy+bn y^2+ax}{2}
                                                        \tag{D.MP4}
\]
has period differences $-ax$ and $by+(b\rho-ab r)/2$.
Here $n,\rho$ are odd. Therefore $\phi$ is integral on every
lattice point. Its vertical slope has odd constant term
$bn/2$, whereas every coefficient of that slope is divisible
by $p$. The line's characters on the rational fixed double
curve $C$ are consequently one for the scalar involution and
zero for $C_p$. The wall degree is odd and a unit modulo $p$,
since its signed value is $-a(\rho-br)$.

At a rank-one site take its actual primitive period $u_0=(u_1,u_2)$
and write $H=\alpha u_0,F=\beta u_0$ with coprime
$\alpha,\beta$. The integral support is
\[
       \phi_1(j)=-\frac{a u_1}{2\alpha}j(j-\alpha).
                                                        \tag{D.MP5}
\]
The extracted period identities give $\alpha\mid a u_1$;
either $\alpha$ is odd or $a u_1/\alpha$ is even. Hence (D.MP5)
is integral. Its two differences are precisely
$-a u_1j$ and $b u_2j+(b\rho-ab r)/2$. The elliptic
zero-period combination has zero exponent in the degenerating
coordinate. These are extensions of both actual factors, with
their given marking, over the primitive toric fans.

For the scalar involution there is exactly one further fixed
curve $E_2$, at an $I_2$ site. Its short period $(2,ar)$ is
primitive and (D.MP5) becomes $-a j(j-1)$, with even slopes.
For $C_p$ there is exactly one further fixed curve $E_p$,
at an $I_p$ site. Its short period $(p,ar)$ is primitive and
the doubled support becomes $-ap j(j-1)$, with slopes
divisible by $p$. The corresponding normal weight is
\[
                         w=(ar)^{-1}\pmod p.             \tag{D.MP6}
\]
There are no other fixed strata for either subgroup. These
claims use the actual primitive two-dimensional period columns,
not merely the Kodaira types. They are checked for every row
in the retained certificate. The original additive caps have
free scalar action. In each simply connected finite cap the
actual gamma numerator is primitive modulo its full order;
applying gamma to an equality with a scalar translate forces
that affine power to be the identity. Thus these caps add no
fixed points even when their orders are divisible by $2$ or $p$.

\emph{The two complete primary invariants.}
The mod-two generator $v$ of $H^2(X;\mathbb F_2)$ has a Borel
lift whose characters on $C,E_2$ are respectively one and zero.
To justify a simultaneous choice, the based scalar action on
$W$ gives $\pi_1(W_{C_2})=\pi_1(W)\times C_2$. Its unique
order-two subgroup is the one seen at every fixed point.
The difference from the just-constructed line, being ordinary
zero on $W$, comes from group cohomology. Its restrictions
differ only by the same base class, and possibly a mixed
degree-one term on $E_2$. One base adjustment therefore gives
exactly the restrictions (D.C216). The normal Euler class on
both curves is $\tau^4$, because both normal signs are negative
and their total Chern degree is even. The same Thom-localization
calculation as in D.2.35 yields
\[
                         \langle v^3,[X]_2\rangle=1.     \tag{D.MP7}
\]
Consequently the reduction of the spin-bordism class of $X$
to $K(\mathbb Z/2,2)$ is the conic-surgery class, by the full
calculation $\Omega^{\mathrm{Spin}}_6(K(\mathbb Z/2,2))
=\mathbb Z/2$ used in D.2.8.

For the odd reduction take the doubled line (D.OP4); its
equivariant characters on $C,E_p$ are both zero. Normalize
$z\in H^2(X;\mathbb F_p)$ by $z[C]=1$. This is possible
because the wall degree in (D.MP4) is a unit. The torsion-free
group $\pi_1(W)$ again makes the base adjustment common to
both fixed curves. Thus the restrictions have exactly the
form (D.OP8). The weights are $(1,-1)$ and $(w,-w)$.
Localization in degree two gives $z[E_p]=-w^2$, the ratio
of the equivariant $p_1$ to the normal Euler class is $-2$,
and the odd-degree mixed term squares to zero. Therefore
\[
 \langle z^3,[X]\rangle=0,\qquad
 \langle z\,p_1(X),[X]\rangle=-2(1-w^2).                 \tag{D.MP8}
\]
D.1.10 identifies the entire target
$H_6(K(\mathbb Z/p,2);\mathbb Z)=\mathbb Z/p$ with this cube
evaluation. Hence the odd fundamental-class invariant is zero,
and the refined Pontryagin coefficient in the stated generator
is $(w^2-1)/2$. Theorem~5 of \cite[p.796]{manZhubr75} says
that these two invariants inject the entire spin-bordism group;
it applies to an arbitrary map to $K(\mathbb Z/p,2)$, so no
isomorphism $H_2(X)\cong\mathbb Z/p$ is being assumed here.
Up to an automorphism its class is exactly that of
$Z_{p,0}$ if $w^2=1$, and $Z_{p,1}$ otherwise.

\emph{Why these primary computations identify the whole manifold.}
We supply the splitting, rather than assuming that ordinary
characteristic classes exclude a mixed refinement. Put
$K_2=K(\mathbb Z/2,2)$ and $K_p=K(\mathbb Z/p,2)$.
All their positive integral homology is respectively two-primary
and $p$-primary. One direct verification uses the path-loop
fibration: after inverting the relevant prime, the positive
homology of the fibre $K(\mathbb Z/q,1)$ vanishes by the
cyclic-group resolution. Its Serre sequence then identifies
the base homology with that of the contractible path space.
Thus every positive integral homology element of the base is
killed by a power of $q$, as claimed.

All tensor and Tor products between positive homology of
$K_2$ and $K_p$ vanish. The integral reduced Kunneth sequence
therefore makes $K_2\wedge K_p$ acyclic, and equivalently
the inclusion
\[
                K_2\vee K_p\longrightarrow K_2\times K_p
                                                        \tag{D.MP9}
\]
an integral homology equivalence. Both spaces are simply
connected CW complexes. The relative Hurewicz theorem,
applied to the first possible nonzero relative homotopy group,
and then Whitehead's theorem make (D.MP9) a homotopy equivalence.
Since $\Omega^{\mathrm{Spin}}_6(\mathrm{pt})=0$, the wedge
axiom for reduced spin bordism consequently gives
\[
 \Omega^{\mathrm{Spin}}_6(K(\mathbb Z/(2p),2))
 \cong\Omega^{\mathrm{Spin}}_6(K_2)
       \oplus\Omega^{\mathrm{Spin}}_6(K_p).               \tag{D.MP10}
\]
Here the comparison to the product is the actual Chinese
remainder isomorphism on $H_2$, and the two components are
its two projections. Its automorphism group is the product
of the corresponding primary automorphism groups.

The connected sum on the right of (D.MP1) maps to this product
by collapsing its separating sphere and using the canonical
classifying maps of its two summands. Its $H_2$ map is the
direct-sum isomorphism. A single oriented spin one-handle
bordism compares the connected sum with the disjoint union,
so its two bordism components are precisely the two just
computed for $X$. Both manifolds are simply connected, spin,
have this same $H_2$, and have $b_3=0$. Zhubr's Theorem~3
\cite[p.795]{manZhubr75} now gives the required oriented
diffeomorphism. In particular the possible two-primary
refinement has been determined by (D.MP7), not discarded.

Finally the actual Tate weights are $w=4$ modulo five for
the three $\epsilon=0$ rows, $w=3$ modulo five for the three
$\epsilon=1$ rows, and $w=4$ modulo seven for all four
remaining rows. This gives exactly (D.MP2). Distinct choices
of $\epsilon$ have distinct complete primary bordism
invariants, and both occur among the constructed members.
\end{proof}

The exact certificate \path{code/manifolds_mixed_prime_AH.py}
retains all twelve patterns and ten rows, reconstructs their
Tate markings, and verifies the primitive angular lattices,
$736$ semicharacter substitutions, every period difference,
integrality and scalar character, both unique fixed elliptic
curves, and all odd normal weights. It imposes no finite
multiplicity or twist cutoff. The topological prime splitting
is proved above; it is not inferred from those finite checks.
